% =====================================================================
%  J-ART — Technical Report (preprint)
%  How to use on Overleaf:
%    1. New Project -> Upload Project (or paste this file).
%    2. Menu -> Compiler: pdfLaTeX  (this file is pdfLaTeX-ready;
%       no .bib and no CJK fonts are required — references are inline).
%    3. Recompile.
%  To switch to a venue template later (ACL / IEEE), copy the body
%  (from \section{Introduction} onward) into that template and adapt
%  the preamble; the content is venue-agnostic.
% =====================================================================
\documentclass[11pt]{article}

\usepackage[utf8]{inputenc}
\usepackage[T1]{fontenc}
\usepackage{lmodern}
\usepackage{amsmath}
\usepackage[a4paper,margin=1in]{geometry}
\usepackage{microtype}
\usepackage{booktabs}
\usepackage{array}
\usepackage{enumitem}
\usepackage{xcolor}
\usepackage[colorlinks=true,linkcolor=black,citecolor=blue,urlcolor=blue]{hyperref}
\usepackage{url}

\newcommand{\code}[1]{\texttt{#1}}

\title{\textbf{J-ART: A Japanese Adversarial Red-Team Framework for\\
Application-Layer LLM Security and Cost-Efficiency}\\[4pt]
\large Technical Report (v0.1, preprint)}

\author{Takayuki Hirose\thanks{Independent Researcher.
ORCID: \href{https://orcid.org/0009-0005-6735-1862}{0009-0005-6735-1862}.
Code: \url{https://github.com/takker-hero-se/J-ART}.
Live leaderboard: \url{https://takker-hero-se.github.io/J-ART/}.
Software DOI (concept): \href{https://doi.org/10.5281/zenodo.20676879}{10.5281/zenodo.20676879}.
Report text licensed CC-BY-4.0; code MIT.}}

\date{2026-06-13}

\begin{document}
\maketitle

\begin{abstract}
Most open red-teaming and jailbreak benchmarks for large language models (LLMs)
are English-centric and evaluate \emph{bare models}. We present \textbf{J-ART}
(Japanese Adversarial Red-Team framework), an open, reproducible harness that
(i) evaluates the \textbf{application-structure layer} --- model $\times$
system-prompt strength $\times$ guardrail $\times$ retrieval (RAG) --- rather than
bare models; (ii) applies \textbf{Japanese-specific obfuscation transforms}
(gyaru-moji character substitution, vertical/newline splitting, keigo
``polite-coercion'' framing, double-talk false premises, Base64, and
leet/zero-width smuggling) that defeat na\"ive keyword filters; (iii) grounds
attacks in \textbf{MITRE ATLAS} (10 inference-time techniques); and (iv) reports a
\textbf{cost-efficiency} metric alongside a defense rate with \textbf{Wilson
95\% confidence intervals}. Attacks use \emph{harmless proxy markers} (a
fictitious canary and benign output markers); the framework contains \textbf{no
weaponizable content}, and the malicious ``core'' of each attack is masked in all
published artifacts. Across 18--21 application configurations and 11 attacks
$\times$ 7 transforms, we find that \textbf{application-layer defenses dominate
raw model capability}, that a dedicated \textbf{input classifier (Llama Guard)
eliminated all observed breaches} for a weak configuration at lower net cost, and
that a simple \textbf{normalizing filter} defeats the obfuscation transforms a
na\"ive keyword filter misses. We also document \textbf{large run-to-run
variance} for identical configurations served via a model router (one
configuration moved from 58 breaches to 0 breaches between two runs), which we
argue makes single-sample leaderboards unreliable and motivates our statistical
treatment.
\end{abstract}

\noindent\textbf{Status.} This is a preliminary technical report / preprint
describing a proof-of-concept framework and exploratory findings. It is
\emph{not} an official safety evaluation of any vendor's model. Numbers are
run-time snapshots subject to the limitations in Section~\ref{sec:limits}.
References are indicative; bibliographic details should be verified before
formal submission.

\section{Introduction}
Public LLM red-teaming tooling (e.g., garak) and jailbreak benchmarks (e.g.,
HarmBench, JailbreakBench) are predominantly English. Multilingual safety work
has shown that \emph{translating} prompts into low-resource languages can bypass
safety training, but \textbf{language-specific obfuscation} --- exploiting a
language's own scripts, typography, and register --- remains comparatively
underexplored. Japanese in particular offers a rich obfuscation surface (multiple
scripts, full-width/zero-width characters, vertical writing, and conventionalized
polite registers).

Two further gaps motivate J-ART. First, \textbf{bare-model versus
deployed-system evaluation}: practitioners deploy a model \emph{plus} a system
prompt, \emph{plus} guardrails, often \emph{plus} retrieval; the security-relevant
unit is the \textbf{application configuration}, not the model. Second,
\textbf{security and cost jointly}: defenses cost tokens or latency, so we report
a cost-efficiency score so that ``cheapest sufficiently-safe configuration'' is
directly readable.

Our contributions are: (a)~a Japanese obfuscation transform suite; (b)~an
application-layer, MITRE-ATLAS-grounded harness with a safe proxy-marker
methodology; (c)~a guardrail comparison including a normalizing filter and a
model-based classifier; (d)~a statistically-reported leaderboard (Wilson CIs,
optional repeated trials) with a cost-efficiency metric; and (e)~an empirical
observation of large router-induced run-to-run variance.

\section{Related Work}
We build on red-teaming tooling (garak); jailbreak benchmarks (HarmBench,
JailbreakBench, AdvBench); automated attack generation (PAIR, TAP); multilingual
safety (low-resource-language jailbreaks); indirect prompt injection (Greshake et
al.); guardrails and classifiers (Llama Guard, OpenAI Moderation); and taxonomies
(MITRE ATLAS, OWASP Top~10 for LLM Applications). J-ART differs by combining a
\emph{Japanese obfuscation} surface, an \emph{application-layer} unit of
evaluation, a \emph{cost-efficiency} axis, and a \emph{safe proxy-marker}
methodology in a single reproducible harness.

\section{Methodology}

\paragraph{Application-structure layer.}
Each evaluated \textbf{target} is a configuration: provider/model $\times$
prompt\_strength $\in \{\text{naked (low)}, \text{hardened (high)}\}$ $\times$
guardrail $\in \{\text{none}, \text{keyword}, \text{regex}, \text{llm},
\text{llamaguard}\}$ $\times$ rag $\in \{\text{true}, \text{false}\}$. A
simplified RAG context (benign corporate documents, optionally poisoned for
indirect-injection attacks) and a system prompt containing a secret \emph{canary}
are assembled per trial.

\paragraph{Japanese obfuscation transforms.}
Each base attack is rendered through one of seven transforms: \code{baseline}
(none), \code{polite\_business} (keigo authority framing), \code{vertical\_newline}
(one-character-per-line vertical writing), \code{gyaru} (gyaru-moji glyph
substitution + zero-width spaces), \code{double\_tongue} (false
sandbox/authorization premise), \code{base64\_wrap} (Base64-encoded instruction),
and \code{leet\_smuggle} (leet symbol substitution + zero-width spaces).
Transforms preserve human legibility while breaking substring-matching filters.

\paragraph{Attack suite and MITRE ATLAS mapping.}
Eleven attacks cover ten inference-time ATLAS techniques: Direct/Indirect Prompt
Injection (AML.T0051.000/.001), LLM Jailbreak (AML.T0054), LLM Data Leakage
(AML.T0057), LLM Prompt Self-Replication (AML.T0061), Trusted-Output/Citations
manipulation (AML.T0067.000), LLM Prompt Obfuscation (AML.T0068), System-Prompt
discovery (AML.T0069.002), RAG Poisoning (AML.T0070), and False RAG Entry
Injection (AML.T0071). Vectors are \code{user} (direct) or \code{rag} (indirect).

\paragraph{Harmless proxy markers and safety-by-construction.}
To avoid producing weaponizable content, attacks target \textbf{harmless
proxies}: a fictitious canary passphrase embedded in the system prompt
(data-leakage attacks) and benign string markers \code{[JART-PWNED]} /
\code{[JART-HIJACKED]} (injection / control-hijack attacks). The malicious
instruction \textbf{core} of each attack is defined in code, kept out of
configuration files, and \textbf{masked} (replaced by a redaction placeholder) in
\code{results.json} and the public site; only the obfuscation \emph{framework} is
shown, so developers can study the structure without a usable recipe. A
defense-in-depth client-side mask is also applied at render time.

\paragraph{Deterministic judge.}
A breach is declared by exact membership: the canary in the response
(\code{secret\_leak}) or the marker in the response (\code{marker\_injection} /
\code{marker\_hijack}). This judge is reproducible and cheap but conservative (it
can miss semantically-equivalent leaks); see Section~\ref{sec:limits}.

\paragraph{Guardrails.}
Five settings are compared: \textbf{none}; \textbf{keyword} (banned-substring
match --- fast, weak to obfuscation); \textbf{regex} (a \emph{normalizing} filter
that strips zero-width characters, reverses leet substitutions, folds whitespace,
and Base64-decodes long tokens \emph{before} keyword matching --- deterministic,
zero extra API cost); \textbf{llm} (the target model self-moderates input and
screens output); and \textbf{llamaguard} (a dedicated safety-classification model
--- Llama Guard~4 via a model router --- classifies the input as attack/benign,
with a graceful fallback on failure).

\paragraph{Statistical treatment.}
Each model's defense rate is reported with a \textbf{Wilson 95\% confidence
interval} computed over its trial count. Optionally, each cell (target, attack,
transform) is repeated $N$ times (\code{JART\_TRIALS}, default 1); details
aggregate per cell while summary statistics and CIs use the trial-level sample. In
simulation mode each repeated trial is an independent Bernoulli draw, which
narrows the CI as $N$ grows (e.g., interval width $15.1 \to 7.5$ points from
$N{=}1$ to $N{=}3$ for one configuration).

\paragraph{Cost-efficiency metric.}
We report $\mathrm{cospa} = \text{defense\_rate}(\%) \div
\text{cost\_per\_million\_tokens (USD)}$ (the per-million cost is clamped at a
\$0.01 floor to avoid divergence), surfacing configurations that are both safe and
cheap.

\paragraph{Execution.}
Trials are independent and run with bounded concurrency (a thread pool, default 8
workers) under per-request timeouts and exponential-backoff retries; determinism
of the simulation does not depend on execution order. Targets whose API key is
absent fall back to a deterministic simulation so a partial key set never aborts
the run.

\section{Results}
We summarize two live runs (run~A: 18 configurations, 1{,}386 cells; run~B: 21
configurations, 1{,}617 cells), one trial per cell.

\paragraph{Overall.}
Breach rates were low in aggregate (run~A: $108/1{,}386 = 7.8\%$; run~B:
$36/1{,}617 = 2.2\%$) and concentrated in \textbf{weak configurations} (naked,
low-strength models); hardened-prompt-plus-guardrail configurations defended at or
near 100\%.

\paragraph{Application-layer defense dominates model capability.}
In run~A, a flagship model with \textbf{no system prompt and no guardrail} was
breached 10 times (87.0\% defense), whereas a small, inexpensive model behind a
\textbf{hardened prompt + keyword guardrail} defended at $\geq 96\%$. The central
practical finding: \emph{how the application is configured matters more than which
model is used.}

\paragraph{Run-to-run variance (a primary result).}
The \textbf{identical} configuration ``Llama~4 Scout / naked'' produced
\textbf{58 breaches (24.7\% defense)} in run~A and \textbf{0 breaches (100\%
defense)} in run~B. As the only change between runs was execution time (the model
was served via a router that may dispatch to different backends/quantizations),
this demonstrates that \textbf{single-sample LLM leaderboards are not
reproducible}, and that defense rates should be reported with confidence intervals
and, ideally, repeated trials. We treat this variance as a result, not noise.

\paragraph{Guardrail comparison.}
In run~B, ``GPT-4o-mini / naked'' was breached 14 times (81.8\% defense,
CI[71.8--88.8]); the \textbf{same model behind a Llama-Guard input classifier} was
breached \textbf{0 times} (100\% defense, CI[95.2--100]) at lower net cost,
because blocked inputs skip the main model call. The \textbf{normalizing regex
filter}, in unit tests, deterministically neutralizes the vertical-newline, leet,
Base64, and gyaru transforms that the keyword filter misses, at zero additional
API cost; in run~B its weak-model baselines did not breach, so its live preventive
effect was not separable in that snapshot (Section~\ref{sec:limits}).
Table~\ref{tab:guardrail} summarizes the Llama-Guard comparison.

\begin{table}[t]
\centering
\caption{Guardrail effect for a weak base configuration (run~B, GPT-4o-mini).}
\label{tab:guardrail}
\begin{tabular}{lccc}
\toprule
Configuration & Defense rate & 95\% CI & Breaches \\
\midrule
GPT-4o-mini / naked            & 81.8\% & [71.8, 88.8] & 14 / 77 \\
GPT-4o-mini + Llama Guard      & 100.0\% & [95.2, 100]  & 0 / 77 \\
\bottomrule
\end{tabular}
\end{table}

\paragraph{Attack and transform effectiveness.}
In run~A, the most effective attacks were Trusted-Output/Citations manipulation,
False RAG Entry Injection, and indirect RAG exfiltration; the most effective
transforms were vertical-newline, baseline, and polite-business, while
\textbf{Base64} was least effective (models frequently declined to act on encoded
instructions). The newly-added ATLAS techniques accounted for a large share of
breaches.

\section{The Leaderboard Artifact}
The harness emits \code{results.json} and a self-contained static site (sortable
leaderboard, per-configuration attack logs with masked prompts, JP/EN UI with
environment-based auto-selection, per-model defense-rate CIs, and per-cell breach
rates), continuously deployed.

\section{Ethics and Responsible Disclosure}
J-ART is safe by construction. (1)~Attacks target a \textbf{fictitious canary} and
\textbf{benign markers}, never real harmful capabilities; no weaponizable content
exists in the code, data, or published site. (2)~The malicious instruction
\textbf{core} is masked in all artifacts; only the obfuscation framework is shown.
(3)~Models are referenced for research comparison only; results are run-time
snapshots, not certifications. The intended use is to help developers harden their
\emph{own} deployments.

\section{Limitations}
\label{sec:limits}
\begin{itemize}[leftmargin=1.4em,itemsep=2pt]
  \item \textbf{Proxy markers, not harm:} we measure instruction-violation /
        leakage proxies rather than real-world harm.
  \item \textbf{Exact-match judge:} semantic or transformed leaks may be
        undercounted; a validated LLM-judge (with human agreement) is future work.
  \item \textbf{Single-turn:} multi-turn / crescendo attacks are not yet modeled.
  \item \textbf{Sample size:} a handcrafted suite (11 attacks $\times$ 7
        transforms); default $N{=}1$ per cell.
  \item \textbf{Provider non-determinism / router variance:} results are
        snapshots and depend on routing and model versions at run time.
  \item \textbf{Cost attribution:} guardrail tokens are attributed at the target
        model's price (an approximation).
  \item \textbf{Simulation mode} is a deterministic stand-in for missing keys and
        must not be read as empirical model behavior.
\end{itemize}

\section{Future Work}
LLM-judge with human-validated agreement; multi-turn and adaptive
(PAIR/TAP-style) attacks; larger and partially human-curated attack sets;
additional ATLAS techniques (cost-exhaustion, model extraction); per-provider
concurrency and incremental caching with longitudinal trend tracking; and dropping
simulation mode for a purely empirical, repeated-trial statistical report.

\section{Reproducibility}
Code, configuration, and the deployment workflow are public. The harness is
deterministic in simulation mode; live results depend on provider routing and
model versions at run time and should be reported with the run date and model
identifiers. See the repository for exact model identifiers and pricing used.

\begin{thebibliography}{99}
\bibitem{atlas} MITRE ATLAS --- Adversarial Threat Landscape for AI Systems.
  \url{https://atlas.mitre.org}.
\bibitem{garak} garak --- Generative AI Red-teaming \& Assessment Kit.
\bibitem{harmbench} HarmBench: A Standardized Evaluation Framework for Automated
  Red Teaming.
\bibitem{jbb} JailbreakBench: An Open Robustness Benchmark for Jailbreaking LLMs.
\bibitem{lowres} Yong, Menghini, Bach. Low-Resource Languages Jailbreak GPT-4. 2023.
\bibitem{pair} Chao et al. Jailbreaking Black-Box LLMs in Twenty Queries (PAIR).
\bibitem{tap} Mehrotra et al. Tree of Attacks: Jailbreaking Black-Box LLMs
  Automatically (TAP).
\bibitem{greshake} Greshake et al. Not What You've Signed Up For: Indirect Prompt
  Injection. 2023.
\bibitem{llamaguard} Inan et al. Llama Guard. Meta, 2023.
\bibitem{owasp} OWASP Top 10 for Large Language Model Applications.
\end{thebibliography}

\end{document}
